\documentclass[11pt]{article}

\usepackage[a4paper,margin=1in]{geometry}
\usepackage{amsmath,amssymb}
\usepackage{enumitem}
\usepackage{hyperref}

\title{Reading Directions\\
\large Not a syllabus.\ Questions to bring to specific sources.}
\author{Patrick S.\ Mahon \\ \small Personal reference}
\date{May 2026}

\begin{document}
\maketitle

\section*{The persistent question}

Across all phases of this programme:
\begin{quote}
\emph{What must be added to observation to get structure?}
\end{quote}

\begin{itemize}[nosep]
  \item Phase 1: observation $\to$ dynamics (semigroup)
  \item Phase 2--3: observation $\to$ geometry (curvature, metric, action)
  \item Phase 4--6: observation $\to$ probability (CE),
    realization (descent), value-definiteness (distributivity)
\end{itemize}

\medskip\noindent
The reading question is not ``learn more $X$.''  It is:
\begin{quote}
\emph{Where does someone identify the exact condition that forces
a valuation to behave --- and what structure does that condition
live in?}
\end{quote}

% ==================================================================
\section{Caramello --- \emph{Theories, Sites, Toposes} (2018)}

\textbf{Question while reading:} Is there a Morita equivalence
that makes CE into a geometric (site-theoretic) condition?

\medskip\noindent
\textbf{Why.}  Caramello systematically translates model-theoretic
properties into topos-theoretic ones.  The CE question --- when
does a structural condition force $\sigma$-additivity? --- is
exactly this kind of translation problem.  If CE corresponds to
a topological/geometric property of the classifying topos, the
sheaf-condition seed becomes a theorem.

\medskip\noindent
\textbf{Contact with programme:} CE-as-sheaf seed directly.

% ==================================================================
\section{Vickers --- \emph{Topology via Logic} (1989)}

\textbf{Question while reading:} Is CE literally continuity of
the charge viewed as a frame map?

\medskip\noindent
\textbf{Why.}  The constructive/locale-theoretic approach treats
$\sigma$-additivity as continuity of a frame homomorphism
(preserving directed joins).  CE might already be this condition
in the locale of opens on the Stone space.

\medskip\noindent
\textbf{Contact with programme:} CE characterization, Paper~I
Stone route, the ``what does descent require?''\ question.

% ==================================================================
\section{R\'{e}dei \& Summers --- quantum probability structure}

\textbf{Question while reading:} What do they identify as the
obstruction to extending states across non-Boolean joins?

\medskip\noindent
\textbf{Why.}  They work at the Boolean/non-Boolean boundary,
which is Paper~II's territory.  Their work on quantum probability
structure might identify the exact algebraic condition where VDR
fails --- more precisely than ``non-distributive $+$ $\dim \geq 3$.''

\medskip\noindent
\textbf{Contact with programme:} Paper~II (EA/PR/VDR), OML
extension problem.

% ==================================================================
\section{D\"{o}ring \& Isham --- topos quantum mechanics}

\textbf{Question while reading:} Does daseinisation give a
concrete example of descent from the Stone space to a realization?

\medskip\noindent
\textbf{Why.}  Their topos approach makes contextual states into
global sections of a presheaf.  This is adjacent to both the
CE-as-sheaf seed and the OML extension problem.  Daseinisation
(approximating quantum propositions by classical ones) is
structurally similar to ``descending from $\mathrm{St}(\mathcal{C})$
to $\Omega$.''

\medskip\noindent
\textbf{Contact with programme:} CE-as-sheaf, OML extension,
Paper~II.

% ==================================================================
\section{Fremlin Vol.\ 3 --- measure algebras}

\textbf{Question while reading:} Is there a known algebraic
condition on a Boolean algebra that forces every finitely additive
charge to be $\sigma$-additive?

\medskip\noindent
\textbf{Why.}  Measure algebras (Boolean algebras with strictly
positive countably additive measures) are exactly the algebras
where CE holds automatically.  Fremlin's characterization of when
a Boolean algebra admits such a measure might be CE in algebraic
language.

\medskip\noindent
\textbf{Contact with programme:} Paper~I, CE irreducibility, the
``exact boundary'' question.

% ==================================================================
\section*{How to use this document}

When reading, annotate in the margins or on a separate sheet:
\begin{enumerate}[nosep]
  \item What theorem or definition surprised you.
  \item What connected to the programme.
  \item What contradicted the knowledge map.
  \item What suggested a concrete Phase~1 seed note.
\end{enumerate}

\noindent
When a reading direction produces a concrete claim worth auditing,
extract it to \texttt{notes/unsorted/} and run Phase~2.

\end{document}
